% language=us

\environment context-2023-style

\startdocument
  [title={challenge},
   banner={a bit of fun},
   location={context\enspace {\bf 2023}\enspace meeting}]

% \setuphead
%   [chapter]
%   [after=\blank] % {\blank[line]}] fails with figure left

\setupbackgrounds
  [text]
  [text]
  [state=stop]

\defineoverlay[whatever][\processMPbuffer]

\starttitle[title=The Sundqvist Smiley]

This is how teens use potrace:

\startbuffer
fill
    lmt_potraced [ bytes =
        "001111100
         010000010
         100000001
         101101101
         100000001
         101000101
         100111001
         010000010
         001111100",
         polygon = false,
         size = 1,
    ] xysized (OverlayWidth-5mm,OverlayHeight-5mm)
    withcolor "extracolor"
    withtransparency (1,.3)
    withpen pencircle scaled 0.1 ;
\stopbuffer

\typebuffer[style=\switchtobodyfont[9pt]\ttbf]

\stoptitle

\starttitle[title=An alternate view]

And this is how we modulate on that:

\startbuffer
fill
    lmt_potraced [ bytes =
        "..11111..
         .1.....1.
         1.......1
         1.11.11.1
         1.......1
         1.1...1.1
         1..111..1
         .1.....1.
         ..11111..",
         polygon = false,
         size = 1,
    ] xysized (OverlayWidth-5mm,OverlayHeight-5mm)
    withcolor "extracolor"
    withtransparency (1,.3)
    withpen pencircle scaled 0.1 ;
\stopbuffer

\typebuffer[style=\switchtobodyfont[9pt]\ttbf]

And we can go further, see next page.

\stoptitle

\starttitle[title=Some more detail]

\startbuffer
draw image (
    lmt_startpotraced [ bytes =
        "..11111..
         .1.....1.
         1.......1
         1.22.22.1
         1.......1
         1.3...3.1
         1..333..1
         .1.....1.
         ..11111.."
    ] ;
    fill lmt_potraced [ value = "1", size = 1 ]
        withcolor "extracolor" withtransparency (1,.3)
        withpen pencircle scaled 0.1 ;
    fill lmt_potraced [ value = "3", size = 1 ]
        withcolor "extracolor" withtransparency (1,.5)
        withpen pencircle scaled 0.1 ;
    fill lmt_potraced [ value = "2", size = 0 ]
        withcolor "extracolor" withtransparency (1,.8)
        withpen pencircle scaled 0.1 ;
    lmt_stoppotraced ;
) xysized (OverlayWidth-5mm,OverlayHeight-5mm) ;
\stopbuffer

\typebuffer[style=\switchtobodyfont[9pt]\ttbf]

\stoptitle

\starttitle[title=This year X-mas card (1)]

Here is how we use it ourselves:\high{\copyright 2023}

\startbuffer[1]
\startluacode
local sin, cos = math.sin, math.cos

local N = 100

local function f(x,y)
    local z = sin(x*x) - cos(y*y)
 -- local z = sin(x) - cos(y)
    if z > 0 then
        return '1'
    else
        return '0'
    end
end

potrace.setbitmap("mybitmap", potrace.contourplot(N,N,f))
\stopluacode
\stopbuffer

\startbuffer
path p ; p := lmt_potraced [
    stringname = "mybitmap",
    value      = "1",
    size       = 0,
] ;
draw image (
    draw p withcolor "extracolor" withtransparency (1,.25) ;
  % fill p withcolor "darkred"    withtransparency (1,.3) ;
) xysized (OverlayWidth-5mm,OverlayHeight-5mm) ;
\stopbuffer

\getbuffer[1]

\typebuffer[style=\switchtobodyfont[12pt]\ttbf]

We also turned this into a game we played during the holidays and that later got
properly materialized by Willy. We also made a service so that users could
generate matching games in different complexity.

\stoptitle

\starttitle[title=This year X-mas card (2)]

\typebuffer[1][style=\switchtobodyfont[12pt]\ttbf]

\stoptitle

\stopdocument
